\chapter{二、三的奇次幂：同一结构的另外两个周期}\label{ch:twothree}
\section{迹与行列式先告诉我们该试什么}
对一般三块核心 $(u^k,v^r,w)$，$w=-ku-rv$，平均 $u^{k-r},v^{r-1},w$ 给出整数矩阵
\[C=\begin{pmatrix}-r-1&1\\-n&1\end{pmatrix},
\qquad n=k+r+1.
\]
它满足
\[\operatorname{tr}C=-r,\quad\det C=k,\quad C^2+rC+kI=0.
\]
忽略共同尺度后，控制特征值之比的是
\[\theta=\frac{r^2}{k}.
\]
若该比值为单位根，$C$ 才有有限射影阶。因为 $2+2\cos\phi=\theta$ 是有理数，且单位根给代数整数，有理的 $2\cos\phi$ 只能是 $-2,-1,0,1,2$。在 $r>0$ 且矩阵非标量的本模板内，真正的有限阶恰好对应
\[\theta=1,2,3,
\]
射影阶分别为三、四、六。$\theta=4$ 是非平凡Jordan情形，并不产生有限阶。\src{S21}

平方族给 $\theta=1$。若 $k=ms^2,r=ms$，则 $\theta=m$，所以 $m=2,3$ 正好提供另外两类周期。这一计算指导我们选取模板，但不能把 $\theta>4$ 解释成平均问题不可能可解；它只说明这个 $C$ 不再自行循环。

\section{相同的两步选点，新的库存比例}
以下令 $m=2$ 或三，$s\ge2$，并取
\[k=ms^2,\qquad r=ms,\qquad n=ms^2+ms+1.
\]
有
\[C^4=-k^2I\quad(m=2),\qquad C^6=-k^3I\quad(m=3).
\]
因此实际返回 $C/k$ 可逆。另两个一步返回是
\[
\mathcal A=\begin{pmatrix}1&-1/s\\0&-1/s\end{pmatrix},\qquad
\mathcal B=\frac1k\begin{pmatrix}-r-1&r\\-n&k+r\end{pmatrix}.
\]
由 $C^{-1}(k\mathcal B)=\Sigma$，得到 $\Sigma$；再由 $\Sigma\mathcal A=\Delta(1/s)$ 开始建立倍率资源。\src{S17}

对每个 $1\le d\le r$，先选
\[(u^{k-r+d-1},v^{r-d},w)
\]
平均。保留 $r$ 份新值和一份未选的旧 $u$，再平均其余 $k$ 个位置。两步矩阵的整数部分为
\[
M_d=\begin{pmatrix}
r^2+r-k&-rd\\rn&-(r+k)d
\end{pmatrix}=C^2\Delta(d),
\]
实际尺度为 $k^{-2}$。因为 $C$ 有正向逆，直接得到全部 $\Delta(d)$。此处最大 $d=r$ 仍留一份旧 $u$，最小 $d=1$ 也不超出 $v$ 的库存，所以整个参数区间都是实际可用的。

\section{恢复逆元为什么分成几个小情况}
与平方族一样，初始倍率 $s$ 和 $\Sigma$ 只给带因子 $s-1$ 的上剪切，以及 $R_0=\Z[1/s]$ 上的深层。要消掉固定因子，仍希望得到二的可逆倍率。

先补 $m$ 的逆。如果 $m\mid s$，已有的 $\Delta(s)^{-1}\Delta(s/m)$ 直接给出 $\Delta(1/m)$。如果 $m\nmid s(s-1)$，倍率与旧水平互素，有限商方法适用。剩下 $m\mid s-1$ 时，构造
\[
H_m=\frac1{mk}\Delta(m)C^2\Delta(m)
=\begin{pmatrix}
((m-1)s+1)/(ms)&-1/s\\
n/s&-m(s+1)/s
\end{pmatrix}\in\SL_2(R_0).
\]
左上角为 $R_0$ 中元素，正是因为 $s\equiv1\pmod m$。前篇的深层工具给 $H_m$ 的逆，把这个“夹心”式子解开就能得到 $\Delta(1/m)$。

当 $m=3$、$s$ 为奇数时，还需要另补除二。相应行列式一矩阵为
\[
H_2=\frac1{2k^2}\Delta(2)C^4\Delta(2)
=\begin{pmatrix}
(s+1)/(2s)&-1/s\\n/s&-4-2/s
\end{pmatrix}.
\]
它的分母二同样被 $s+1$ 的偶性消去。于是各个分支都得到 $\Delta(2)^{\pm1}$，重算交换子给 $U(1/2)$。

令 $R=\Z[1/(2ms)]$，最终资源与平方族相同：
\[U(R),\quad L(nR),\quad\text{安全小倍率及其逆}.
\]
下剪切来自 $U(-1)(C/k)\Delta(k)=L(-n/k)$。整个推导对所有指定 $m,s$ 都成立。

\section{入口和小单位问题也能接上}
任意输入的入口仍用见证转移。不过这次有更方便的恒等式：
\[m(k-1)-(r-2)(r+2)=4-m.
\]
结合 $n=k+r+1$ 可得 $\gcd(k-1,n)=1$：$m=3$ 时公因子整除一；$m=2$ 时整除二，而 $n$ 为奇数。因此三角色交换子在每个危险素数处都有效。保护位置数量由 $\omega(n)\le r$ 控制，入口成立。

剩下单位群由 $-1$ 和 $1\le d\le r$ 的安全小整数生成。GRH下的小非剩余界在 $s\ge256$ 已落入 $r=ms$ 以内；$2\le s<256$ 的两族共508个参数已有精确有限生成证书。它们接上通用控制器和终端，得到：
\begin{theorem}[二、三平方倍数族]
在Dirichlet $L$ 函数的GRH下，对 $m\in\{2,3\}$、$s\ge2$，$k=ms^2$ 平均在 $n=ms^2+ms+1$ 个位置上恰以 $G=1$ 判定。
\end{theorem}
小参数证书所覆盖的每个具体维数是无条件的。十二平均十九元、十八平均二十五元、五十平均六十一元都属于这两个族；但这些维数未必是各自最早成功点。

\section{为什么这完成了二、三的所有幂的最早成功点}
当 $k=p^a$，临界尺度为 $S(k)=p^{\lceil a/2\rceil}$。偶指数写成平方，上一讲给 $k+S(k)+1$；奇指数且 $p=2,3$ 时写成 $ps^2$，本讲给同一位置数。结合严格下界，对于 $a\ge2$，在GRH下有
\[M(p^a)=H(p^a)=p^a+p^{\lceil a/2\rceil}+1,
\qquad p=2,3.
\]
素数 $k=3$ 由前篇处理；$k=2$ 是 $M=4,H=5$ 的特殊自动点。

这仍未给出 $N(p^a)$。例如九平均十三元成功，没有顺带证明十四元。一个核心解决一个维数；把后续全部维数连接起来，是另一项独立任务。
